% META: DONE.1

\section{Связь с игровой индустрией}

\subsection{Общие точки соприкосновения CAD и игровой индустрии}

На первый взгляд может показаться, что системы автоматизированного проектирования (САПР), традиционно ассоциирующиеся с машиностроением и промышленным дизайном, имеют мало общего с индустрией разработки видеоигр. Однако за последнее десятилетие эти две области значительно сблизились, а их технологические границы стали всё более прозрачными.

Современные игровые движки, такие как Unreal Engine и Unity, всё чаще используются не только для создания развлекательного контента, но и для промышленной визуализации, симуляции и создания цифровых двойников (digital twins) производственных объектов \cite{sastwin2025}. В свою очередь, CAD-системы - в частности, AutoCAD - активно применяются для создания базовых 3D-моделей игровых сред и объектов \cite{autodeskgaming2017}. На стыке этих областей сформировался целый ряд направлений, в которых разрабатываемый в рамках данной магистерской диссертации метод восстановления истории построения трёхмерных моделей может найти прямое и востребованное применение.

\subsection{Проблема массового производства 3D-моделей для видеоигр}

Разработка современных видеоигр, особенно в жанрах с открытым миром, требует колоссального объёма трёхмерного контента. Как отмечают разработчики из Monolith Soft (создатели серии Xenoblade Chronicles), ещё несколько лет назад для создания игры требовалось от 1 000 до 2 000 моделей вручную, но теперь эта цифра достигла 100 000, что делает ручное создание каждого актива практически невозможным \cite{assetgeneration2025}. «Обычный метод имел ограничения с точки зрения затрат и сроков», - констатируют разработчики, отмечая, что автоматизация части процесса стала необходимостью.

Традиционный пайплайн создания игровых 3D-активов включает следующие этапы:
\begin{enumerate}
    \item Концептуализация и создание 2D-эскизов.
    \item Моделирование в специализированном ПО (Blender, Maya, 3ds Max).
    \item Текстурирование и наложение материалов.
    \item Оптимизация полигональной сетки под требования игрового движка.
    \item Импорт и настройка в движке.
\end{enumerate}

Этот процесс остаётся трудоёмким и слабо автоматизированным, что приводит к значительным временным и финансовым затратам, поэтому здесь методы, разработанные в настоящем исследовании, могут внести существенный вклад.

\subsection{Прямое применение результатов исследования в игровой индустрии}

Разработанный в рамках диссертации модуль «НейроКОМПАС-3D», осуществляющий восстановление истории построения трёхмерных моделей, может быть адаптирован для решения нескольких ключевых задач игровой индустрии.

\subsubsection{Восстановление истории построения форматов игровых ассетов}

Видеоигры хранят свои ассеты в специализированных, зачастую недокументированных форматах для моделей, текстур, анимаций, аудио и уровней. Сообщество моддеров и исследователей игр постоянно занимается обработкой этих форматов, чтобы извлекать, конвертировать и модифицировать игровые модели. Существуют целые коллекции инструментов и документации, посвящённые этой задаче \cite{awesomeGameReversing}.

Методы, разработанные в данной работе - в частности, нейросетевая сегментация геометрии и распознавание формообразующих операций - могут быть адаптированы для:
\begin{itemize}
    \item \textbf{Распознавания структуры игровых моделей}: Игровые модели, извлечённые из проприетарных форматов, часто представляют собой «немую» полигональную сетку без информации о том, как она была создана. Применение предложенного в работе гибридного подхода (сегментация + детерминированное восстановление) позволит автоматически распознавать базовые примитивы, из которых составлена модель, и реконструировать её параметрическую историю.
    \item \textbf{Автоматической параметризации моделей для моддинга}: Создание модов (пользовательских модификаций) для игр часто ограничено отсутствием инструментов для редактирования готовых моделей. Возможность восстановить параметрическое представление по итоговой сетке открыла бы для моддеров возможность легко изменять размеры, формы и другие параметры игровых объектов, не имея исходных файлов в редакторе.
\end{itemize}

\subsubsection{Оптимизация пайплайна CAD-to-Game Engine}

Как отмечается в исследованиях, посвящённых конвертации CAD-моделей в игровые ассеты, существуют серьёзные проблемы при переносе детализированных CAD-моделей в игровые движки: потеря параметрической иерархии, проблемы оптимизации сетки, отсутствие материалов и цветов. Современные решения, такие как плагин RapidPipeline для Unreal Engine \cite{rapidPipeline}, позволяют выполнять очистку, упрощение и оптимизацию CAD-моделей непосредственно внутри редактора.

Однако ни одно из существующих решений не предлагает восстановления истории построения - то есть понимания того, из каких примитивов и операций состояла исходная CAD-модель. Разработанный в данной работе модуль мог бы стать важным дополнением к таким инструментам, позволяя:
\begin{itemize}
    \item Автоматически распознавать формообразующие операции импортированной модели.
    \item Создавать на основе распознанной структуры параметрическое представление, пригодное для дальнейшего редактирования в игровом движке или во внешних инструментах.
    \item Упрощать модели путём последовательного удаления распознанных операций, что особенно полезно для создания уровней детализации (LOD) - важнейшей задачи оптимизации производительности игр \cite{lodWikipedia}.
\end{itemize}

\subsubsection{Применение в генеративных и процедурных пайплайнах}

Процедурная генерация контента (Procedural Content Generation, PCG) становится стандартом де-факто для создания больших открытых миров \cite{pcgTechniques}. Современные исследования в этой области направлены на автоматическое создание интерактивных 3D-миров из текстовых описаний с использованием LLM и диффузионных моделей \cite{llmWorlds}.

В этом контексте разработанный метод может быть использован для:
\begin{itemize}
    \item \textbf{Анализа существующих моделей для обучения генеративных систем}: Размеченная история построения может служить обучающими данными для нейросетей, генерирующих параметрические последовательности команд.
    \item \textbf{Семантической сегментации игровых моделей}: Как отмечается в научной литературе, семантика 3D-моделей играет ключевую роль в играх и симуляциях, а автоматическая сегментация и аннотирование больших наборов моделей - активно исследуемая задача \cite{scitepressSegmentation}.
\end{itemize}

\subsection{Примеры сценариев использования}

Рассмотрим подробнее, в каких типах игр и на каких этапах разработки могут найти применение результаты диссертации.

\subsubsection{Игры с процедурной генерацией контента}

В играх с процедурной генерацией (например, в жанрах roguelike, sandbox, open-world survival) часто требуется автоматическое создание тысяч вариаций объектов - от растений и построек до транспортных средств. Использование предложенного метода для восстановления истории построения позволит:
\begin{itemize}
    \item Создавать параметрические шаблоны объектов на основе анализа существующих моделей.
    \item Автоматически варьировать параметры распознанных операций для генерации новых вариаций.
    \item Обеспечивать согласованность генерируемых объектов (например, все фаски на модели имеют согласованные углы).
\end{itemize}

\subsubsection{Строительные симуляторы и игры с крафтингом}

В играх, где игрок может строить и модифицировать объекты (таких как Minecraft, Satisfactory, Space Engineers), возможности для редактирования часто ограничены сеточной структурой. Методы восстановления могли бы позволить:
\begin{itemize}
    \item Распознавать сложные формы, построенные игроком, и предлагать их параметрическое представление.
    \item Упрощать сложные конструкции путём выделения повторяющихся элементов (например, массива одинаковых отверстий).
    \item Экспортировать игровые постройки в CAD-форматы для дальнейшего использования в промышленных приложениях (например, для 3D-печати).
\end{itemize}

\subsubsection{Обучение и тренажёры с использованием игровых движков}

Игровые движки всё активнее применяются для создания обучающих симуляторов и тренажёров в промышленности \cite{ieeeTrainer}. В таких системах важно не только визуализировать CAD-модели, но и давать пользователю возможность их редактировать и модифицировать. Разработанный модуль позволил бы:
\begin{itemize}
    \item Автоматически создавать параметрические модели для обучающих сценариев на основе готовых STEP-файлов.
    \item Демонстрировать студентам и инженерам логику построения моделей путём пошагового удаления операций.
    \item Обеспечивать обратную связь при обучении: показывать, какая операция была пропущена или выполнена неверно.
\end{itemize}

\subsubsection{Игры на основе реальных локаций и объектов}

Современные игры всё чаще используют фотограмметрию и 3D-сканирование для переноса реальных объектов и локаций в игровую среду \cite{realityCaptureTutorial}. Однако полученные таким образом сетки имеют крайне высокую плотность треугольников (до 50 млн треугольников) и нерегулярную топологию \cite{instalod}, что делает их непригодными для прямого использования в реальном времени.

Разработанный в диссертации подход, основанный на семантической сегментации сетки с помощью MeshCNN, после доработки для предсказания отдельных частей больших моделей, может быть применён для:
\begin{itemize}
    \item Распознавания на отсканированной сетке отдельных объектов и их составных частей.
    \item Восстановления параметрического представления сканированного объекта (например, распознавание того, что данное отверстие является сквозным выдавливанием, а данное ребро - скруглением).
    \item Создания оптимизированных LOD-моделей с сохранением семантической информации.
\end{itemize}

\subsection{Примеры из индустрии и существующие аналоги}

На рынке уже существуют решения, демонстрирующие востребованность подобных технологий:

\begin{itemize}
    \item \textbf{Kaedim}: облачная CAD-платформа, использующая AI и ML для создания игровых ассетов, генерирует модели из четырёхугольников (quad-based) с водонепроницаемой топологией, оптимизированные для игровой разработки \cite{kaedim}.
    \item \textbf{RapidPipeline}: плагин для Unreal Engine, позволяющий импортировать и оптимизировать CAD-модели непосредственно в редакторе, выполняя декацимацию, ремешинг, удаление невидимой геометрии и генерацию текстурных атласов \cite{rapidPipeline}.
    \item \textbf{RealityCapture и Agisoft Metashape}: инструменты фотограмметрии, активно применяемые в игровой индустрии для создания фотореалистичных моделей реальных объектов \cite{ieeePhotogrammetry}.
\end{itemize}

Ни одно из перечисленных решений, однако, не предлагает функциональности восстановления параметрической истории построения по итоговой геометрии - та задача, которая решается в данной диссертации.

\subsection{Потенциальные направления развития для игровой индустрии}

На основе проведённого анализа можно сформулировать следующие направления, в которых разработанный метод мог бы быть адаптирован и развит для нужд игровой индустрии:

\begin{enumerate}
    \item \textbf{Адаптация под игровые форматы}: Расширение поддержки входных форматов с STEP/STL на популярные игровые форматы (FBX, OBJ, glTF, проприетарные форматы игровых движков).
    \item \textbf{Распознавание игровой семантики}: Обучение нейросетей распознаванию не только CAD-операций, но и игровых примитивов (персонажи, оружие, транспортные средства, элементы окружения).
    \item \textbf{Интеграция с игровыми движками}: Разработка плагинов для Unity и Unreal Engine, позволяющих выполнять восстановление истории построения моделей непосредственно внутри редактора.
    \item \textbf{Генерация LOD с сохранением семантики}: Использование результатов сегментации для создания интеллектуальных уровней детализации, где разные части модели могут упрощаться с разной степенью в зависимости от их семантической важности (например, лицо персонажа должно сохранять детализацию дольше, чем элементы одежды).
    \item \textbf{Обучение на игровых датасетах}: Сбор и разметка датасетов игровых моделей для обучения специализированных нейросетей, учитывающих особенности игровой геометрии (низкополигональные модели, особенности LOD, требования к анимации).
\end{enumerate}

\subsection{Заключение}

Таким образом, разработанный в рамках данной магистерской диссертации метод восстановления истории построения трёхмерных моделей, несмотря на свою ориентацию на CAD-системы и машиностроительные приложения, обладает значительным потенциалом для применения в индустрии разработки видеоигр. Общие технологические основы (работа с трёхмерными сетками, семантическая сегментация, нейросетевой анализ геометрии), а также конкретные потребности игровой индустрии - автоматизация создания активов, обратный инжиниринг проприетарных форматов, оптимизация пайплайна CAD-to-Engine - создают прочную основу для дальнейшей адаптации и развития разработанных решений в контексте игровой разработки.

